\assignment{Vector manipulation and matrices}{08/10/2025}
\begin{qst}{Algebra}{x=5}
Definire $ x=5 $ e $ y=2 $ e calcolare i seguenti:
\begin{enumerate}
    \item $ u=x+y $, $ v=xy $, $ \displaystyle w=\frac{x}{y} $, $ z=w^3 $;
    \item $ \displaystyle s=\frac{xy^2}{(x-y)} $, $\displaystyle p=\frac{3x}{2y} $, $\displaystyle r=\frac{3xy}{2} $, $\displaystyle t=\frac{x^5}{(x^5-1)} $. 
\end{enumerate}
\end{qst}   
\begin{exs*}\end{exs*}
\begin{matlab}
x=5;
y=2;

u=x+y
v=x*y
w=x/y
z=w^3
\end{matlab}
\begin{matlabconsole}
>> lab01

u =
     7

v =
    10

w =
    2.5000

z =
   15.6250
\end{matlabconsole}

\begin{exs*}\end{exs*}
\begin{matlab}
x=5;
y=2;

s= (x*y^2) / (x-y)
p= (3*x) / (2*y)
r= (3*x*y) / 2
t= (x^5) / (x^5-1)
\end{matlab}
\begin{matlabconsole}
>> lab01

s =
    6.6667

p =
    3.7500

r =
    15

t =
    1.0003
\end{matlabconsole}

\begin{qst}{Vectors and matrices}{vct}
Definire il vettore $ x=[0:0.1:1] $, digitare i seguenti comandi Matlab e comprenderne il significato:
\begin{enumerate}
    \item $ x([1\text{ }3]) $;
    \item $ x([1:3]) $ (latex spazia per indentazione);
    \item $ x([1 \text{ } 2 \text{ } 3]) $;
    \item $ x([3: \text{end}]) $.
\end{enumerate}
\end{qst}
\begin{exs*}\end{exs*}

\begin{qst}{Matrices}{matr}
Definire la matrice:
\[ 
  A=
  \begin{bmatrix}
    1 & 2 & 3 & 4 \\
    5 & 6 & 7 & 8 \\
    9 & 10 & 11 & 12
  \end{bmatrix}
.\] 
Comprendere il significato dei seguenti comandi:
\begin{enumerate}
    \item \verb|size(A)|;
    \item \verb|B=A.*A|;
    \item \verb|B=A*A|;
    \item \verb|B=A'*A|;
    \item \verb|B=A*A'|.
\end{enumerate}
\end{qst}

\begin{qst}{diag}{diag}
Utilizzare il comando diag (digitare help diag per vederne l'utilizzo) di Matlab per definire
la matrice tridiagonale B di dimensione $ 10 \times 10 $  così definita:
\begin{itemize}
    \item Gli elementi della diagonale principale sono tutti uguali a 5;
    \item Gli elementi della codiagonale inferiore sono tutti uguali a -1;
    \item Gli elementi della codiagonale superiore sono tutti uguali a 3;
    \item Gli elementi appartenenti all'intersezione delle colonne 6 e 9 e delle righe 5 e 8 sono uguali a 2.
\end{itemize}
\end{qst}

\begin{qst}{Matrices}{matr2}
Utilizzare i comandi \verb|zeros|, \verb|ones| e \verb|eye| di Matlab per definire le seguenti matrici 
(digitare \verb|help zeros|, \verb|help ones| e \verb|help eye| per vederne l'utilizzo).

\[ 
  W=
  \begin{bmatrix}
  0 & 0 & 0 & 1 & 1 & 1 \\
  0 & 0 & 0 & 0 & 0 & 0 \\
  0 & 0 & 0 & 0 & 0 & 0
  \end{bmatrix},
\] 
\[ 
  Z=
  \begin{bmatrix}
  1 & 0 & 0 & 0 & 1 & 1 \\
  0 & 1 & 0 & 0 & 1 & 1 \\
  0 & 0 & 1 & 0 & 1 & 1
  \end{bmatrix}
.\] 
\end{qst}